\section{代码与数据开源说明}

本研究坚持可重复科学研究准则，论文所涉及的全部数据处理脚本、集合调整卡尔曼滤波同化模块、交叉熵轨迹规划器、深度学习时空世界模型对比代码以及行前假想预演元控制器的完整源代码与实验配置文件，均已整理并开源于代码托管平台：\TBD{GitHub仓库地址}。

实验所使用的海洋卫星遥感大数据均取自国际公开科学数据中心，具体获取途径与数据集标识如下：
\begin{enumerate}
  \item \textbf{东海多源融合叶绿素a浓度遥感数据集（2018—2026年）：} 来源于美国国家海洋和大气管理局海岸观测网（NOAA CoastWatch）ERDDAP数据服务器，融合了可见光红外成像辐射仪（VIIRS）与哨兵三号海洋陆地颜色仪（OLCI）观测并经经验正交函数插值重构（DINEOF）生成，空间分辨率为2千米，时间分辨率为逐日，遵循CC0公共领域开放许可，数据集唯一标识符为\texttt{noaacwNPPN20S3ASCIDINEOF2kmDaily}。
  \item \textbf{全球多尺度超分辨率海表温度数据集（2018—2026年）：} 来源于美国喷气推进实验室（JPL）物理海洋学分布式活动存档中心ERDDAP服务器，版本为MUR v4.1，原始空间分辨率为$0.01^\circ$，时间分辨率为逐日，数据集唯一标识符为\texttt{jplMURSST41}。
\end{enumerate}

\section{人工智能工具使用说明}

在本课程论文的研究与撰写过程中，作者严格遵循浙江大学关于学术诚信与人工智能工具规范使用的相关规定，采用人机协同模式开展工作，各环节的具体分工与责任界定说明如下：
\begin{enumerate}
  \item \textbf{研究构思、算法设计与实验执行：} 论文的核心科学问题提出（具身采样中的思考—行动权衡与过思考偏差剖析）、东海赤潮遥感大数据清洗与422个任务窗口构建、搜索—留出双集合解耦架构设计、行前假想预演元控制器算法实现以及全部数值仿真实验与统计检验，均在作者主导下配合人工智能研究助手（Claude）共同完成；所有实验数据、统计指标与图表结果均源自真实运行日志。
  \item \textbf{文献调研与核查：} 论文引用的25篇参考文献均经过出版商原始记录或预印本库逐一核对，确保作者、题名、期刊/会议名称、年份及数字对象唯一标识符（DOI）真实准确，未引入任何未经核实的虚假文献。
  \item \textbf{论文初稿撰写与文字润色：} 在给定完整实验事实清单、数学公式定义与图表说明的严格约束下，使用大语言模型（Gemini）辅助生成中文学术初稿的LaTeX文本与中英文摘要，并由Claude与作者逐句核查文中每一项事实陈述、实验数值、公式符号及文献引用的一致性，最终由作者通读审定并对全文学术质量与事实准确性承担全部责任。
\end{enumerate}
